% Generated draft: requires mathematical and source review.
\documentclass[11pt]{article}
\usepackage[margin=25mm]{geometry}
\usepackage{fontspec}
\setmainfont{Latin Modern Roman}
\usepackage{amsmath,amssymb,mathtools,unicode-math}
\setmathfont{Latin Modern Math}
\usepackage{enumitem,needspace,xurl,hyperref}
\setlength{\parindent}{0pt}
\setlength{\parskip}{5pt}
\setlength{\emergencystretch}{3em}
\newcommand{\N}{\mathbb{N}}
\newcommand{\Z}{\mathbb{Z}}
\newcommand{\Q}{\mathbb{Q}}
\newcommand{\R}{\mathbb{R}}
\newcommand{\C}{\mathbb{C}}
\newcommand{\F}{\mathbb{F}}
\newcommand{\A}{\mathbb{A}}
\newcommand{\PP}{\mathbb{P}}
\newcommand{\E}{\mathbb{E}}
\newcommand{\eps}{\varepsilon}
\newcommand{\dd}{\,\mathrm{d}}
\begin{document}
\section*{UnsolvedMath ID: 30002316}
\textbf{Source problem code:} catalogue entry
\subsection{Definitions and mathematical statement}
Let \(k\) be a field. An algebraic \(k\)-torus is an algebraic group \(T\) over \(k\) for which \(T_{k_s}\simeq (\mathbb{G}_{m,k_s})^d\) for some integer \(d\geq 0\), where \(k_s\) is a separable closure of \(k\); its dimension is \(\dim T=d\). The torus \(T\) is anisotropic over \(k\) if it has no nontrivial \(k\)-split subtorus. A field extension \(L/k\) splits \(T\) if \(T_L\simeq (\mathbb{G}_{m,L})^d\).

The supplied question concerns anisotropic algebraic \(k\)-tori admitting a finite Galois splitting extension \(L/k\) whose degree is a prime power. For a prime \(p\) and an integer \(n\geq 1\), the proposed equality is
\[
[L:k]=p^n \quad\Longrightarrow\quad \operatorname{cdim}(T)=\dim T.
\]
The question asks whether this equality holds for every such torus and splitting extension. The prime may vary with the instance, and the source imposes no minimality condition on the splitting extension.

Here \(\operatorname{cdim}(T)\) retains the source's notation for the canonical dimension of the algebraic torus. The supplied record does not define this invariant. In particular, it does not specify a construction using torsors or another convention, nor whether an ordinary or prime-local invariant is intended. This definition must be established from the original source before the equality is a fully specified mathematical assertion; this draft does not identify it with the canonical dimension of the underlying variety.

No base field or characteristic restriction is supplied. Treating the question as universally quantified over fields \(k\) is therefore an unresolved reading of the record, not a verified source hypothesis. The record also does not say whether degree \(1=p^0\) is counted as a prime power. If it is, that boundary case is additionally part of the question; an anisotropic torus split over \(k\) itself is necessarily the trivial torus. No positive-dimension hypothesis is stated.
\subsection{Short English statement}
The question asks whether every anisotropic algebraic torus that becomes split over a finite Galois extension of prime-power degree has canonical dimension equal to its dimension. The supplied record leaves the definition of canonical dimension and any restrictions on the base field unspecified, so these require mathematical review.
\subsection{Sources}
\begin{itemize}
\item Local UnsolvedMath registry: supplied record for permanent ID \(30002316\), title \emph{Canonical Dimension of Prime-Power-Split Anisotropic Tori}, source problem code \texttt{catalogue entry}, in the Oberwolfach Reports Open Problems collection. The record gives proposed year \(2013\) and names no proposer.
\item Supplied source citation: \emph{Algebraic Groups (2014)}, Oberwolfach Reports, DOI \texttt{10.4171/owr/2013/17}, pp.~1025--1085. Supplied source URL: \url{https://doi.org/10.4171/owr/2013/17}. The record describes this material as originating in mathematical workshop open-problem discussions. The bibliographic year is retained exactly as supplied, despite the \texttt{2013} component of the DOI.
\item Supplied literature source, labelled ``Oberwolfach algebraic-groups report'': \url{https://doi.org/10.4171/owr/2013/17}.
\item Literature assessment checked on \texttt{2026-08-21}: the supplied curation labels the problem open and reports that the equality remains unresolved in the reviewed primary literature, with no proof or counterexample found. This is supplied dated curation, not an independently verified literature assessment. The background expressly requests independent verification of time-sensitive status and cited claims.
\item The supplied statement status is \texttt{corrected\_verified}. Its accompanying verification note says that mathematical content was unchanged while notation, typography, and wording were normalized. These are registry assertions; this draft does not independently verify them or the cited report.
\end{itemize}
\end{document}
